\section{Method}
\label{sec:method}
Our system follows a progressive design. Part~1 establishes what can be achieved with spatial resampling, a shallow CNN, and naive temporal aggregation. Part~2 replaces these weak components with pretrained restoration models that either exploit temporal propagation or optimize perceptual realism. Part~3 then studies a low-cost extension that explicitly balances the two behaviors. All methods use the same synthetic degradation protocol on GT datasets: an HR frame is bicubic-downsampled by a factor of 4 to produce the LR input, and the output is resized to the original HR resolution before metric calculation.

\subsection{Part 1: Hand-crafted and Shallow CNN Baselines}
For synthetic datasets with HR frames, we generate LR inputs by bicubic downsampling by a factor of 4. We then reconstruct HR frames with four baselines. This part is important because it reveals two common failure modes: interpolation is stable but blurry, while simple temporal fusion can be harmful when motion is ignored.

\textbf{Bicubic and Lanczos interpolation.}
Bicubic interpolation is used as the simplest spatial resampling baseline. Lanczos interpolation uses a sinc-windowed kernel and often preserves sharper edges than bicubic, although it can introduce ringing.

\textbf{SRCNN.}
We implement a three-layer SRCNN model following the classic design~\cite{dong2015srcnn}. The LR frame is first upsampled by bicubic interpolation to the target HR size, then refined by convolutional layers with kernel sizes $9$, $5$, and $5$. The first layer extracts patch-level features, the second maps them to a compact non-linear representation, and the last layer reconstructs RGB output. SRCNN is frame-independent and therefore does not model temporal consistency.

\textbf{Temporal averaging.}
For each bicubic-upsampled frame $I_t$, we apply a weighted temporal average:
\begin{equation}
\hat{I}_t = 0.25 I_{t-1} + 0.5 I_t + 0.25 I_{t+1},
\end{equation}
with boundary frames clamped to the nearest available index. We then apply a light unsharp mask using a Gaussian-blurred version of the averaged frame. This baseline intentionally does not perform motion compensation, making it useful for demonstrating why alignment is important. If neighboring frames are misaligned by object or camera motion, the average can reduce detail instead of improving it.

\subsection{Part 2: Pretrained SOTA Reproduction}
\textbf{Real-ESRGAN.}
Real-ESRGAN is applied frame by frame using the official x4plus weights. It uses an RRDB-style generator and adversarial/perceptual training to address the regression-to-the-mean effect that often makes distortion-oriented SR blurry. In our implementation, large frames are processed with tiling to reduce VRAM usage. Real-ESRGAN is useful for visual sharpness, but because it is image-based and optimized for perceptual realism, its output may deviate from GT pixels and may not be temporally consistent.

\textbf{BasicVSR and BasicVSR++.}
BasicVSR processes an LR sequence as a tensor of shape $(B,T,C,H,W)$ and uses bidirectional propagation with optical-flow alignment. Information is propagated forward and backward through time, allowing the model to aggregate details from multiple neighboring frames. BasicVSR++ improves this with second-order grid propagation and stronger alignment, which better reuses information from non-adjacent frames. We reproduce both models on mandatory sample data. For the full REDS validation benchmark, we report BasicVSR++ as the stronger VSR representative to control computational cost while preserving the main SOTA comparison.

\subsection{Part 3: Adaptive Hybrid Refinement}
Part~3 targets the trade-off between distortion fidelity and perceptual sharpness. BasicVSR++ provides a stable GT-aligned backbone, while Real-ESRGAN provides sharper perceptual details. We fuse them using:
\begin{equation}
H_t = (1 - M_t)B_t + M_tR_t,
\end{equation}
where $B_t$ is the BasicVSR++ frame, $R_t$ is the Real-ESRGAN frame, and $M_t$ is an adaptive mask clipped to a maximum weight of 0.65.

The mask is computed from three cues. First, edge strength encourages Real-ESRGAN contribution in textured regions. Second, disagreement between $B_t$ and $R_t$ suppresses the mask when Real-ESRGAN deviates strongly from the stable backbone. Third, temporal frame difference in the BasicVSR++ sequence suppresses Real-ESRGAN in high-motion regions where hallucinated details are more likely to flicker. The mask is smoothed with a Gaussian filter. This method does not require additional training and is computationally cheap.

More concretely, we compute three normalized maps: a Sobel edge magnitude map $E_t$, an inter-method disagreement map $D_t=|R_t-B_t|$, and a temporal motion proxy $T_t$ from the difference between neighboring BasicVSR++ frames. The final mask has the form
\begin{equation}
M_t = \mathrm{clip}\left(w_0 + 0.35E_t - 0.45D_t - 0.25T_t, 0, w_{\max}\right),
\end{equation}
where $w_0=0.18$ and $w_{\max}=0.65$ in our experiments. This design keeps BasicVSR++ dominant while allowing Real-ESRGAN to contribute mainly in relatively reliable textured areas.

\begin{figure*}[t]
  \centering
  \includegraphics[width=0.85\textwidth]{pipeline_flowchart.png}
  \caption{Overall experimental pipeline. HR frames are synthetically degraded to x4 LR inputs, processed by Part~1 baselines, Part~2 pretrained methods, and Part~3 adaptive fusion, then evaluated with pixel, perceptual, temporal, and qualitative metrics.}
  \label{fig:pipeline}
\end{figure*}
